\documentclass[../main.tex]{subfiles}
\begin{document}
%==========================================TIÊU ĐỀ TẬP========================================
\begin{table}[h]
    \begin{adjustwidth}{-1cm}{}
        \begin{tabular}{>{\centering\arraybackslash}p{6cm}|>{\raggedright\arraybackslash}p{12.5cm}}
            \multirow{5}{6cm}
            % Cột bên trái: ÉPISODE và số tập
            {\\[-40pt]
                \flaregothic\fontsize{20pt}{20pt}\selectfont \centering
                \color{eptype}HISTOIRE\color{black}\\
                \freshbot\fontsize{72pt}{72pt}\selectfont
                \color{epnum}2
                \\[18pt]
                % \flaregothic\fontsize{14pt}{14pt}\selectfont
                % \color{epattr}BIENVENUE, \uline{\color{Banpire} Mel} !
                % \flaregothic\fontsize{18pt}{18pt}\selectfont
                % \color{epattr}ÉPISODE PERDU
                \color{black}
            }
             &
            % Dòng thứ nhất: tên tập tiếng Nhật
            {\fotskip\fontsize{18pt}{18pt}\selectfont アロエちゃんの\ruby{寮}{りょう}\ruby{生}{せい}\ruby{活}{かつ}
            } \\[6pt]
            % Dòng thứ hai: tên tập tiếng Anh
             & {\cafeteria\fontsize{18pt}{18pt}\selectfont A Succubi's Daily Life}                \\[4pt]
            % Dòng thứ ba: tên tập tiếng Việt
             & {\vnmsans\fontsize{18pt}{18pt}\selectfont Cuộc sống thường nhật của Aloe}          \\[8pt]
            % Dòng thứ tư: Nhân vật xuất hiện
             & {\brian{\fontsize{14pt}{14pt}\selectfont Track used: Dark Ages - Main Theme}
               }
            \\[8pt]
        \end{tabular}
    \end{adjustwidth}
\end{table}
%=========================================TRUYỆN CHÍNH========================================
\color{black}

\pov{Hikawa Kuon}

Tôi vẫn chưa thể tin dược những gì đã xảy ra trong ngày định mệnh đó.

Một cảm giác gì đó khó có thể diễn tả bằng lời. Tôi không biết liệu mình nên cảm thấy vui mừng vì cuối cùng đã có người ở chung, hay lo lắng vì trách nhiệm trên vai ngày càng nặng thêm.

Mới chỉ vài ngày trước thôi, trong nhà chỉ có gia đình tôi, nhà bác Hijaki-san và bà nội tôi.

Nhưng kể từ ngày định mệnh hôm đó, chúng tôi đã có người đầu tiên sống chung dưới mái nhà của gia đình Hikawa chúng tôi.

Mano Aloe-chan.

Một cô gái trẻ với mái tóc hồng và vàng, đôi mắt xanh lục lấp lánh và hai chiếc sừng nhỏ nhưng không đều trên đầu - dấu hiệu của một succubus.

Nhưng trái với hình tượng quyến rũ của loài quỷ hút sinh khí trong truyền thuyết, Aloe-chan giống như một đứa trẻ ngổ ngáo hơn. Lúc nào cũng tự tin thái quá, nhưng cũng dễ lúng túng nếu bị dồn vào đường cùng.

Em ấy từng là một idol. Đã từng có một giấc mơ. Và đã từng mất đi tất cả trong chớp mắt.

Giờ đây, em ấy đang bắt đầu lại cuộc sống, thắp nén lại giấc mơ ấy trong chính gia đình tôi.

Tôi vẫn nhớ cuộc họp gia đình vào bữa tối ngày Aloe-chan chính thức dọn đến.

Và tôi phải nói là\dots\ tất cả mọi người đều không có ý kiến gì về việc này.

\begin{center}
    \uline{Hồi tưởng\\Ngày hôm đó.}
\end{center}

Mọi người quây quần bên bàn ăn, nơi gia đình hai bên trong gia đình tôi hiếm khi có một bữa tối ghép như thế này. Tôi nhìn quanh một lượt.

Bố mẹ tôi, bác Hijaki, bà nội và Ryuuhou đều đã có mặt đầy đủ. Chỉ có Aloe-chan là thành viên mới trong buổi họp mặt này.

``\vie{\hl{Cô bé đó là ai vậy?}}'' ``\vie{\hl{Cái chị có cặp sừng kia là ai? Bạn gái của Onii-san à?}}'' tôi có thể nghe những lời như vậy từ mẹ tôi và em gái tôi bàn bạc từ phía sau.

Tôi hắng giọng: ``\vie{Trước khi ăn, con có chuyện muốn thông báo.}'' Dẫu biết rằng bố tôi và bác Hijaki-san đã biết, tôi vẫn sẽ nói điểu này cho những người còn lại để họ không bị bỡ ngỡ.

Mọi ánh mắt đổ dồn về phía tôi. Tôi hít một hơi, rồi nói rõ ràng: ``Từ bây giờ, Mano Aloe-chan sẽ cùng nhau sống chung với gia đình ta một thời gian.''

Vào đúng khoảnh khắc đó, em gái tôi khẽ trố mắt ngạc nhiên. Việc các thành viên của \hll\ tới nhà tôi kể từ sau dịp Tết hôm đó, em tôi không ngạc nhiên, nhưng việc một cô gái xa lạ tự nhiên sống chung với gia đình tôi lại là một câu chuyện khác.

``Aloe-onee-chan sẽ sống với chúng ta sao?'' liếc nhìn cô nàng Succubus đang vẫy tay một cách e thẹn. ``E-Em cũng không ngờ tới đấy!''

Tôi gật đầu, giữ giọng điệu bình thản: ``Nói ra thì đúng là rất khó tin, nhưng\dots\ đúng là như vậy.''

Mẹ tôi lên tiếng, nhìn khuôn mặt lo sợ của Aloe-chan mà không giấu nổi sự tò mò, cất giọng trêu: ``Con nhặt được em ấy ở đâu vậy?''

Tôi thở dài, giải thích: ``Chẳng phải nhặt đâu mẹ ạ. Tự em ấy muốn ở chung với chúng ta đấy.''

Cô nàng Succubus vội lên tiếng: ``K-Không phải đâu cô chú ạ! YAGOO cho phép cháu ở đây!''

``Như lời em ấy nói đấy,'' tôi chìa hai tay về phía em ấy, khẽ gật đầu như đồng tình với những lời em ấy nói. ``Chuyện nói ra thì khá là dài, nhưng mà tóm gọn lại thì\dots\ em ấy không có nơi nào khác để đi, nên anh quyết định để em ấy ở đây.''

``Chính xác hơn, ông sếp của nó bảo vậy,'' bác Hijaki-san nói, rót thêm một chén rượu.

``Khoan khoan!'' Ryuuhou giơ tay làm động tác tạm dừng, trừng mắt nhìn tôi với đầy sự cẩn trọng. ``Anh có chắc chị ấy không có ý đồ gì không? Trông xảo quyệt ghê\dots''

Ngay lập tức, Aloe gõ mạnh lên bàn, mắt lườm em gái tôi: ``Này nhóc! Chị là succubus đấy! Nếu chị thực sự muốn ra tay với anh trai em, thì anh ấy đã bị hút cạn sinh khí từ lâu rồi!''

Tôi cầm bát cơm lên, vẫn giữ sự điềm tĩnh vốn có: ``Để xem.''

Ryuuhou há hốc mồm nhìn tôi, rồi quay sang nhìn Aloe, rồi lại nhìn tôi.

``Anh bình tĩnh quá mức đấy!'' em tôi giơ tay chỉ về phía tôi và Aloe-chan liên tục. ``Một idol đã giải nghệ, một succubus lang thang, sống ngay tại nhà mình? Nghe vô lý thế nào ấy\dots''

``Anh đã bảo là nó sẽ rất khó tin mà,'' tôi nhún vai đáp lại.

``Rất hợp lý là đằng khác!'' Aloe-chan cười đầy kiêu ngạo. ``Chị hát hay, đẹp gái, cá tính. Ở đâu mà có một onee-san xinh đẹp thế này miễn phí hả? Em nên biết quý trọng đi hen!''

``Chưa biết chị tài cán thế nào chứ chỉ việc chị là một succubus cũng đủ em cảm thấy cạnh gác rồi\dots'' Ryuuhou đáp lại, tỏ vẻ dè dặt.

Cùng lúc này, mẹ tôi, người đang chuẩn bị bê bát canh vào bếp để lấy thêm, mỉm cười châm chọc: ``Quý hóa quá! Con trai mẹ cuối cùng cũng mang một cô gái về nhà!''

Tôi cảm thấy một cơn lạnh sống lưng. Dường như lúc nào mỗi lần tôi đưa một (hoặc một vài) thành viên từ công ty lên chơi là y như rằng bà ấy sẽ đùa như vậy.

``Mẹ, đừng đùa kiểu đó nữa. Nó bắt đầu phát nhàm rồi đấy,'' tôi cố gắng giữ bình tĩnh, nhưng ánh mắt của mẹ tôi cho thấy bà đang rất thích thú.

Ngay lập tức, vị khách bất chợt bật chế độ lém lỉnh, khoanh tay gật gù như thể vừa tìm ra chân lý: ``Aha\~{} Cháu biết ngay là cô sẽ ủng hộ mà!''

Mẹ tôi bật cười khẽ, đặt tay lên má, tỏ vẻ thích thú. ``Aloe-chan, cháu cứ ở lại nhà này bao lâu tùy thích.''

Ánh mắt bà ấy sau đó liếc về phía tôi, mỉm cười: ``Kuon à, mẹ rất muốn xem con trai mẹ sẽ xoay sở như thế nào.''

Tôi thở dài. Tại sao tôi cảm thấy mình bị phản bội thế này? Bà ấy thật sự người con trai đã chống chọi với đủ sự hỗn tạp ở công ty từ hai năm vừa qua sao?

``Chú hy vọng cháu sẽ hòa hợp thật tốt với mọi người tại nơi này,'' bố tôi lên tiếng, cạch chén với bác Hijaki, mắt liếc nhìn em ấy với sự trìu mến.

``Cảm ơn chú đã tin tưởng cháu! Cháu hứa sẽ chăm sóc Kuon-kun thật tốt!'' Aloe-chan nhanh nhảu đáp lại, cười hớn hở.

Em gái tôi nhìn Aloe-chan, rồi cười trừ với tôi: ``Vất vả cho anh rồi, Nii-san. Có vẻ như công việc \hll\ của anh vẫn sẽ không buông tha kể cả khi ở nhà nhỉ.''

Tôi chỉ có thể thở dài: ``Ừ, anh cũng đã tính trước rồi. Hy vọng là em ấy sẽ không làm loạn\dots''

\begin{center}
    \uline{Kết thúc hồi tưởng.}
\end{center}

Dường như tất cả mọi người đều đồng ý với việc sẽ để Aloe-chan ở lại đây sinh sống với chúng tôi.

Cuộc họp gia đình hôm đó kết thúc trong một bầu không khí vừa yên bình nhưng cũng không giấu nổi sự kỳ lạ. Dù vẫn còn nhiều thắc mắc về sự xuất hiện của Aloe-chan, mọi người dần chấp nhận rằng từ nay gia đình Hikawa sẽ có thêm một thành viên mới.

Những ngày sau đó, cuộc sống thường nhật của chúng tôi bắt đầu có những thay đổi dù nhỏ nhưng đáng kể. Sự xuất hiện của một cô nàng succubus lắm chiêu như Aloe-chan khiến ngôi nhà không còn yên tĩnh như trước, đồng thời kéo theo hàng loạt tình huống dở khóc dở cười.

Với những câu chuyện này, tôi sẽ để Mano Aloe-chan kể lại cho các bạn nghe.

\pov{Mano Aloe}

\stpt{-Histoire 1-}
\stcap{Làm bếp}

Sáng hôm đó, tôi dậy sớm hơn thường lệ. Không phải vì có việc gì gấp gáp, mà vì tôi có một ý tưởng hay ho muốn thực hiện: tôi sẽ nấu bữa trưa cho cả nhà gia đình Kuon-kun!

Từ khi chuyển vào sống ở đây, tôi đã được đối xử rất tốt. Cô Kaien lúc nào cũng dịu dàng, Ryuuhou dù hơi đề phòng nhưng vẫn khá dễ thương, còn Kuon-kun\dots\ thì vẫn là Kuon-kun, vừa chu đáo nhưng cũng đáng ghét theo một cách nào đó. Dù sao đi nữa, tôi cũng muốn thể hiện sự biết ơn theo cách của mình!

Tôi bước xuống bếp, tự tin như một người đầu bếp dày dạn kinh nghiệm. Và rồi, tôi nhìn thấy Ryuuhou đã ở đó từ lúc nào.

``Chị đã bao giờ nấu ăn chưa, Aloe-onee-chan?'' - con bé hỏi, ánh mắt vừa tò mò vừa có chút nghi ngại.

Tôi chống nạnh, tự hào vỗ ngực: ``Nấu nhiều chứ! Chứ không thì làm sao mà chị tồn tại ở nơi này được?''

Hình như chú Ryujin với Kuon-kun chưa nói gì với nhóc đó câu chuyện của tôi à?

Ryuuhou không nói gì, chỉ khẽ gật đầu rồi lấy tạp dề ra. ``Được rồi, em sẽ giúp chị. Em không mấy khi vào bếp đâu, nhưng em biết làm vài món cơ bản.''

``Được thôi!''

\ct

Chúng tôi cùng nhau chuẩn bị nguyên liệu. Tôi quyết định làm một bữa trưa đầy đủ, gồm cơm, cá nướng, trứng cuộn và một ít rau xào. Không phải là hộp cơm nhỏ gọn thường thấy mà là hẳn một mâm cơm đàng hoàng.

Tôi nhanh chóng xử lý cá, thêm một chút gia vị trước khi đem nướng. Trong khi đó, Ryuuhou phụ trách đánh trứng và chuẩn bị rau. Dù mới làm việc cùng nhau lần đầu, hai chị em chúng tôi phối hợp khá ăn ý.

``Ồ? Mùi thơm thế\~{}'' một giọng nói nhẹ nhàng nhưng tỏ ý châm chọc của Kaien vang lên từ phía sau. Tôi quay đầu lại, nhìn thấy bà ấy đang mỉm cười nhìn chúng tôi.

``Tự dưng cháu lại trổ tài làm bếp à?'' bà ấy hỏi, nhìn những món ăn mà hai chị em chúng tôi chuẩn bị.

``Cháu muốn bày tỏ lòng biết ơn thôi ạ!'' tôi nhanh nhẹn đáp, trong khi cô nhóc phụ bếp của tôi cũng trả lời: ``Chị ấy muốn trả ơn nhà mình nên mới vậy ạ.''

``Tốt ghê. Để xem cháu thể hiện thế nào nào. Có hơn được cô không đây?''

``Cháu đảm bảo với cô là hơn ạ!'' tôi đáp lại một cách đắc thắng.

Tôi gật đầu, tập trung hoàn thành nốt bữa ăn. Với kinh nghiệm làm bếp của tôi (chủ yếu là để làm cho qua ngày) và sự giúp sức của Ryuuhou, cuối cùng thì hai chị em chúng đã xong mâm cơm.

Chúng tôi bày tất cả chúng lên mâm, rồi mang xuống phòng của cậu Tổng, nơi mà anh ấy vẫn đang làm việc vào những ngày giáp Tết.

\ct

Chúng tôi ngồi vào bàn, đối diện nhau. ``Mời cả nhà ăn cơm!'' chúng tôi cùng nhau đồng thanh.

Kaien gắp một miếng cá nướng, nhâm nhi một chút, rồi cười nhẹ: ``Cháu khá đảm đang đấy. Cá nướng vừa phải, vị cũng khá vừa miệng.''

Tôi nở một nụ cười đắc thắng. Mẹ của nhà anh ấy khen kìa!

Ryuuhou cũng gật đầu, gắp thêm một miếng trứng cuộn do chính tôi chuẩn bị: ``Em thích phần trứng cuộn này, mềm hơn em tưởng.''

Tôi lại càng tỏ vẻ tự hào, khi có được hai cái gật đầu từ cô em gái và mẹ của anh ấy.

Nhưng quan trọng là, Kuon-kun sẽ nghĩ món ăn tôi làm như thế nào?

Anh ấy ăn một cách chậm rãi, nuốt xuống. Rồi, anh thản nhiên nhận xét: ``Cũng ăn được. Rau giòn, món trứng vừa đủ độ mặn và độ ngọt, cá thì cũng được ướp mịn.''

Nghe những gì anh ấy nhận xét, tôi mừng rỡ, giơ nắm đấm lên cao. ``Tuyệt quá!'' Cuối cùng thì Kuon-kun cũng đã-

``Nhưng nếu so với bố thì\dots''

Tôi sững lại.

Khoan đã.

Cái gì cơ?

``N-Nếu so với bố anh thì\dots?'' tôi lập tức hỏi lại, nhíu mày đầy nguy hiểm.

Tôi thực sự không muốn nghe cụm từ ``nhưng'' ấy một chút nào. Không lẽ anh ấy lại chơi bài vạch lá tìm sâu như hôm nọ sao?

Anh Tổng chớp mắt, không hề có chút cảnh giác nào mà đáp tỉnh bơ: ``\dots thì vẫn còn một khoảng khá xa đấy. Cùng lắm là cũng ngang Ryuuhou thôi.''

*BỐP!*

Tôi đập tay xuống bàn. Anh ấy vừa mở miệng khen câu trước, vậy mà sau đó lại đem so sánh tôi với Ryujin sao!?

``Nè! Ý của anh là đồ của em nấu dở đúng không!?''

Kuon-kun thở dài, vẫn không hề nao núng: ``Anh nói là nó ăn được, chứ có bảo dở đâu.''

``Tại sao đồ của em lại chỉ `ăn được' chứ!?''

Anh ấy đặt bát xuống, rồi chỉ vào từng món mà tôi đã cất công làm: ``Cá chưa rút hết xương, rau thi chỉ có vị chưa có mùi tỏi, trứng thì không có mùi thơm.''

Tôi trừng mắt nhìn cậu ta, rồi quay sang Kaien cầu cứu. Nhưng cô ấy chỉ cười nhẹ, như thể đang thưởng thức một cảnh tượng thú vị. ``Cháu không cần phải bận tâm đâu. Cứ nhắc đến đồ ăn là nó coi bố nó là nhất đấy.''

``Dù gì thì món của chị cũng ngon hơn món em làm rồi mà,'' Ryuuhou thì có vẻ hơi lo lắng, vội vàng can thiệp. ``Aloe-onee-chan bình tĩnh nào.''

Tôi khoanh tay, nghiến răng suy nghĩ. Anh ấy dám chê trình nấu ăn của tôi? Không đời nào!

``Được rồi. Em sẽ cho anh thấy chị có thể làm tốt hơn nữa!''

Tôi hạ quyết tâm. Nếu muốn thu phục cái miệng độc địa của Kuon-kun, tôi cần luyện tập thêm. Và cách tốt nhất chính là tìm đến người mà Kuon vừa nhắc tới - bố anh ấy, chú Ryujin!

\stpt{-Histoire 2-}
\stcap{Đấu vật tay}

Một ngày khác, tôi chẳng có gì để làm. Dạo quanh nhà một lúc, tôi nhìn thấy Kuon-kun đang ngồi đọc sách trên ghế gỗ, trông bình thản một cách đáng ghét.

Một ý tưởng lóe lên trong đầu tôi.

Tôi chạy đến trước mặt cậu ta, đặt hai tay lên hông, rồi hắng giọng thật khí thế: ``Kuon! Đấu tay với ta đi! Ta đây sẽ giành chiến thắng cho mà xem!''

Anh ấy từ từ ngước mắt lên khỏi trang sách, nhìn tôi một giây, rồi đảo mắt xuống cánh tay nhỏ nhắn của tôi. ``Với cánh tay ấy á?''

Tôi đập tay lên bàn, gật đầu cái rụp: ``Chắc chắn! Em có sức mạnh của một succubus! Nếu anh nghĩ anh có cơ hội thắng thì lầm to rồi!''

Ngay bên cạnh đó, Hijaki buông tờ báo xuống, nhướng mày đầy thích thú: ``Ồ? Có trò vui à? Được rồi, tao sẽ làm trọng tài.''

Tôi nở nụ cười đắc thắng, trèo lên ghế ngồi đối diện ahh Tổng, đặt khuỷu tay lên bàn. Anh ấy trông vẫn có vẻ hơi lưỡng lự, nhưng rồi cũng đặt tay mình lên, nắm lấy tay tôi.

``\textbf{Chuẩn bị\dots\ bắt đầu!}''

\dots

\dots

\dots

\textbf{*RẦM!*}

Cánh tay tôi bị Kuon-kun đè bẹp xuống bàn mà anh ất thậm chí còn chưa thay đổi biểu cảm.

Tôi tròn mắt, bàng hoàng đến mức không thốt lên được lời nào trong vài giây.

``K-\textbf{Không thể nào!!} N-Nhưng, em là succubus mà!? Sao anh có thể-!?'' tôi lắp bắp, nhìn cánh tay của tôi rồi ngước lên anh ấy.

Trọng tài trận đấu cười khẩy: ``Mày yếu quá đấy, nhóc. Đừng có giở giọng kiêu ngạo ở đây.''

Tôi trừng mắt nhìn ông ta, phân trần: ``Cái này gọi là bất công! Anh ấy chơi ăn gian!''

Kuon-kun chỉ nhún vai: ``Làm gì có chuyện anh gian lận. Anh cũng là dạng người lười chảy thây đấy thôi.''

``Có mà sức mày tập trung hết vào mồm miệng nhiều rồi nên ít cơ đấy,'' Hijaki châm chọc.

Ta nghe anh ấy như vậy càng tức hơn. Anh ấy cứ ra vẻ mình hơn đây à!?

Thôi, ta đây chả thèm đấu với anh nữa! Chuyển sang đối tượng khác!

\dots

\dots Có rồi.

Nhìn sang Hijaki, ta bỗng lóe lên một ý tưởng. Ông ta dù gì cũng có tuổi rồi, tóc bạc phơ ra kia kìa.

Chắc chắn thể lực của ông ta cũng chẳng còn như trước. Nếu thắng ông ta, ít nhất có thể gỡ gạc lại danh dự!

Ta chỉ tay vào ông ấy, mạnh dạn tuyên bố: ``Vậy bác đấu với cháu đi! Bác còn chê cháu yếu đúng không? Để xem ai thắng ai!''

Hijaki bật cười, đặt điện thoại xuống rồi vặn cổ tay một chút rồi thản nhiên ngồi xuống đối diện với ta.

``Mày muốn thì chiều.''

Ta hít một hơi sâu, nắm lấy tay ông ta, lần này quyết tâm hơn bao giờ hết.

``\textbf{Bắt đầu!}''

\dots

\dots

\dots

\textbf{*RẦM!*}

Cánh tay ta một lần nữa bị đè bẹp xuống bàn, nhưng lần này ta có cảm giác như xương cổ tay mình suýt gãy.

``Ái\dots!''

Ta kêu lên một tiếng đầy đau đớn, ôm tay nhảy dựng lên: ``BÁC CÓ CẦN DÙNG LỰC MẠNH ĐẾN VẬY KHÔNG!?''

Đáp lại, Hijaki chỉ hừ một tiếng, nhấc điện thoại lên đọc tiếp như chẳng có gì xảy ra.

Ta ngồi bệt xuống ghế, xoa cổ tay nhăn nhó đầy uất ức. Chẳng lẽ\dots\ sức mạnh succubus của ta thực sự yếu đến thế sao?

Không thể chấp nhận chuyện này được! Ta phải tìm ra một cách để thắng!

Và rồi, một ý tưởng khác lóe lên trong đầu ta. Nếu ta không thể đọ họ bằng sức mạnh, sao không thử đọ bằng sức bền nhỉ?

Ta nở một nụ cười nham hiểm, chống tay lên bàn: ``Vậy ta đây đòi trận tái đấu! Nhưng lần này, chúng ta sẽ thi xem ai có trụ được trên bàn nhậu lâu hơn!''

Aloe ta vốn dĩ có tửu lượng khá tốt. Nếu đấu về khoản này, chắc chắn ta sẽ cầm chắc chiến thắng!

Ngay khi ta vừa dứt lời, một giọng nói trầm thấp vang lên từ phía sau:

``Ồ, vậy thì để chú tham gia thử nhé?''

Cả phòng đột nhiên im lặng.

Ta quay đầu lại và thấy Ryujin đứng khoanh tay trước cửa, nở nụ cười trìu mến nhưng cảm tưởng lại thấy đầy nguy hiểm.

Ta cảm thấy hơi rùng mình, nhưng kiêu hãnh của một con succubus như ta không cho phép mình lùi bước.

Đã ném lao thì phải theo lao thôi!

``Được thôi! Đấu thì đấu!''

\begin{center}
    \ul{Mười phút sau.}
\end{center}

\dots Chuyện quái quỷ gì thế này\dots?

Ta nằm dài trên bàn, mặt đỏ bừng, đầu óc quay cuồng. Trước mặt là một hàng ly trống rỗng.

Ryujin thì vẫn bình thản như chưa có gì xảy ra, chậm rãi rót thêm một ly rượu cho mình.

Bên cạnh, Kuon-kun cầm một tách trà, chậm rãi nhấp một ngụm như thể toàn bộ chuyện này chẳng liên quan gì đến mình.

T-Ta yếu ớt lẩm bẩm: ``\dots Aloe ta đây\dots\ không phục\dots ''

Hijaki nghe vậy phá lên cười: ``Đừng có mà thách đấu linh tinh nữa, nhóc.''

Ngay sau câu nói như xiên thẳng vào lòng tự trọng đó, ta dần chìm vào giấc ngủ từ lúc nào.

\dots Thật đáng ghét\dots

\stpt{-Histoire 3-}
\stcap{Karaoke}

Ngày 29 tháng Chạp, nếu tính theo Âm lịch, hoặc ít nhất là ta đoán vậy.

Tết là dịp để tụ họp, ăn uống và vui chơi, nhưng với ta - một cựu idol (ta thực sự không thích nói điều này) thì không gì tuyệt vời hơn một buổi karaoke gia đình!

``Nào, ai sẽ mở màn đây?'' cô Kaien tươi cười hỏi mọi người.

Ta lập tức giơ tay như một đứa trẻ háo hức: ``Cháu! Cháu! Cháu sẽ hát!''

Ngước nhìn xuống Kuon-kun, ta thấy anh ấy đang chỉ chăm chăm vào cái điện thoại, trông có vẻ không hào hứng gì.

\dots Thôi thì, kệ anh ấy vậy.

Trong khi đó, chú Ryujin nhấp một ngụm bia, mắt nhìn ta đầy cân nhắc. ``Hừm, để xem Aloe-chan có cải thiện được gì sau lần bọn chú góp ý hôm đó không.''

Ta không hề nao núng với những lời thách thức đó. Lần trước, họ nói giọng ta bị ``cục một vó'', rồi chê là ``quá phấn khích''\dots\ Hôm nay, ta sẽ chứng minh rằng ta đã cải thiện sau nhiều ngày tự luyện hát!

Ta chọn một bài hát sôi động, cầm micro mà ta hay mang theo với khí thế của một chiến binh. Nhạc vang lên, và ta cất giọng.

Bầu không khí thay đổi ngay lập tức. Giọng hát của ta, trong trẻo nhưng lại ẩn chứa một nét quyến rũ đặc trưng của một succubus, tràn ngập không gian.

Ryuuhou mắt mở to: ``\dots Chị ấy hát tốt thật.'' Ta nghe thoáng thoảng như vậy.

Kaien gật gù, ra vẻ trầm ngâm. ``Hmmm\~{} Nếu theo nghề idol, chắc con bé này cũng thành công đấy.''

Ryujin dừng động tác uống bia lại một chút, chú ý tới giọng ca của ta: ``Tốt hơn kha khá so với lần đó rồi đấy.''

Kuon đơn giản chốt lại một câu: ``Đồng ý. Giá mà em ấy không bị dính dáng cái vụ việc kia thì tốt biết mấy.''

Bài hát kết thúc.

Ta cười rạng rỡ, đặt tay lên hông, đầy tự hào. ``Thế nào, thế nào? Cháu có xứng danh là idol không!?''

Kaien cười nhẹ: ``Rất tốt, nhưng thiếu một chút gì đó\dots''

Lại ``thiếu một chút gì đó.'' Thực sự, ta bắt đầu thấy khó chịu với cái gia đình này rồi đấy!

``Lần này là ta lại bỏ sót cái gì nữa?'' mặt ta bày tỏ sự không vui.

Ryuuhou lập tức nói thẳng, nghiêm túc phán: ``Thiếu sự nghiêm túc!''

Hể? Ta không hiểu ta đang sai ở đâu đấy!

\dots Khoan, cả nhà đều đồng tình với ý kiến đó ư!? Cái gì chứ!?

Ta bĩu môi, cảm thấy bị phản bội: ``Ơ kìa, hát karaoke thôi mà!''

Kuon lắc đầu, chống tay lên cằm: ``Hát karaoke thì đúng, nhưng ít nhất cũng phải hát tử tế để mọi người cùng thẩm định chứ!''

Vậy là từ một buổi hát chơi, ta lại rơi vào một buổi huấn luyện bất đắc dĩ. Kaien phân tích cách lên xuống giọng, Ryujin chỉnh lại cách giữ hơi, còn Kuon thì liên tục góp ý về cách thể hiện cảm xúc.

Ta khoanh tay, cau có: ``Đây là karaoke, chứ có phải audition (buổi thử giọng) đâu\dots''

Ryujin chỉ nhếch mép cười khẩy, còn Kuon thở dài, nhưng trong mắt lại ánh lên sự hài lòng.

Không phải hài lòng vì ta bị ``bầm dập'', mà bởi vì họ được mọi người quân tâm.

Nhưng quan tâm kiểu vạch lá tìm sâu thế này thì ta không ưa một chút nào.

Dù sao thì, được hát cùng mọi người cũng không tệ. Còn hơn là hát một mình cho mình ta nghe.

\stpt{-Histoire 4-}
\stcap{Hái lộc và cầu nguyện}

Buổi chiều mùng Một, tiết trời lạnh nhưng không quá rét. Ta đang lười biếng nằm dài trên sô-pha vì không có gì để làm thì bất chợt Kaien xuống phòng ta vỗ tay một cái, khiến tôi và Ryuuhou giật mình.

``Nào, mọi người chuẩn bị đi hái lộc đầu năm thôi!'' cô ấy thông báo với một giọng điệu hớn hở.

Ta chớp mắt, ngồi thẳng dậy: "Hái lộc? Đó là gì thế ạ?"

Cô nhóc của anh Tổng quay sang tôi, giải thích với vẻ hào hứng: ``Là một phong tục đầu năm! Mình sẽ đi chùa, cầu nguyện, rồi hái một nhánh cây để lấy may.''

Ta xoa cằm suy nghĩ, rồi bất chợt hỏi một câu đầy vu vơ: ``Hmmm, vậy có giống lấy năng lượng từ linh hồn người khác không?''

Cả phòng im lặng trong ba giây.

Rồi Kuon-kun - người cũng vừa xuống tới phòng - nhìn ta với ánh mắt bất lực: ``Tuyệt đối không.''

Vậy thì nó trông thế nào nhỉ?

\ct

Sau khi chuẩn bị xong, cả gia đình Hikawa bọn ta cùng nhau đi đến ngôi chùa lớn nhất trong vùng - theo ta được biết là chùa Nomu. Dọc đường, ta nhận ra không khí Tết ở đây thật nhộn nhịp. Dù không giống hoàn toàn với lễ hội ở Nhật Bản hay thế giới của xứ Quỷ giới của ta, nhưng vẫn mang nét gì đó rất đặc biệt.

Ngôi chùa này nằm trên một ngọn đồi nhỏ ở ngoại ô thành phố. Xung quanh nó rất thoáng đãng và toát lên một vẻ trang nghiêm tới đáng sợ. Từ xa, tôi đã thấy rất nhiều người xếp hàng trước khu vực cầu nguyện.

``Trước tiên, mình sẽ vào lễ và cầu nguyện đã,'' - cô Kaien nói, dẫn cả nhóm chúng ta đến khu vực chính điện.

Ta quan sát mọi người cẩn thận, rồi làm theo họ: bỏ tiền lẻ vào hòm công đức, vỗ tay hai cái, chắp tay lại cầu nguyện.

Không biết tại sao ta phải làm vậy. Ta là quỷ mà, phải làm những trò gì đó ranh ma chứ!

\dots

\dots Ranh ma?

Phải rồi! Để xem nơi này ta có thể làm được gì\dots

Để ta thử xem\dots\ năng lực succubus có tác dụng với thần thánh không nhỉ?

Nghĩ là làm, ta nheo mắt, dồn một chút tà khí đặc trưng của mình ra ngoài. Không phải để làm gì xấu đâu nhé! Chỉ là một con quỷ 15 tuổi này \dots\ nổi hứng tò mò chút thôi!

Tuy nhiên, ngay lúc đó\dots

``Cô gái kia, đừng có đứng đó phát ra tà khí nữa.'' Một giọng nói vang lên ngay sau lưng ta, trầm và đầy quyền uy.

Ta giật mình quay lại, và suýt nữa nhảy dựng lên.

Là sư thầy của ngôi chùa! Ông ấy nhìn ta với ánh mắt nghiêm nghị, như thể vừa bắt quả tang ta làm điều cấm kỵ!

``Ối! Cháu không cố ý!'' ta luống cuống giơ tay lên biện hộ.

Nhưng rồi, đã quá muộn.

Kuon-kun và Ryuuhou đứng kế bên, nhìn ta chằm chằm đầy khó chịu. Cô nhóc thì thở dài, còn anh Tổng chỉ lẳng lặng lắc đầu.

``Ngoài đời không như mơ đâu, Aloe-onee-chan ạ!'' - cô nhóc lo xa đó kết luận, khoanh tay như một giáo viên đang phê bình học sinh nghịch dại - là ta.

Ta chỉ có cười trừ, nhanh chóng thu hồi tà khí và lùi sang một bên, cố gắng tỏ ra ngoan ngoãn như chưa có chuyện gì xảy ra.

Chậc, mọi thứ lại không giống như trong truyện tranh hay phim hoạt ảnh rồi. Bực ghê.

\ct

Sau khi cầu nguyện xong, cả nhóm đi ra phía sân chùa để hái lộc.

Ta nhìn mọi người chọn những cành cây nhỏ với vẻ tò mò nhưng rất cẩn trọng. Cô Kaien chọn một nhánh có lá xanh tươi, Kuon-kun thì lấy một nhánh đơn giản. Ryuuhou chọn cành nào đó trông hơi nhỏ, nhưng cô bé lại có vẻ rất hài lòng.

Đến lượt ta.

Nhìn xung quanh một vòng, ta chọn một cành cây trông có vẻ ngầu nhất: một cành có lá dài, hơi uốn lượn, có chút ánh sáng phản chiếu trên bề mặt.

``Hehe, nhìn thế này mới hợp với khí chất succubus chứ!'' Ta vừa lẩm bẩm vừa huơ huơ nhánh cây trên tay.

Nhưng đúng lúc đó, Kuon nhướng mày nhìn ta: ``Aloe-chan, em có chắc em hiểu hái lộc là gì không?''

Ta dừng lại, chớp mắt. ``Hả? Chứ không phải cứ chọn cành nào đẹp là được à?''

Ryuuhou phì cười, giơ nhánh cây của mình lên: ``Không hẳn đâu chị. Quan trọng là tâm niệm khi hái lộc ấy. Người ta tin rằng những nhánh cây này mang đến may mắn cho cả năm, nên không phải cứ lấy đại là được đâu.''

Ta nhìn nhánh cây của mình, rồi lại nhìn nhánh của Ryuuhou. Cái của ta đúng là trông ngầu thật, nhưng nhìn kỹ lại thì\dots\ thật sự chẳng có gì đặc biệt cả.

Ta chớp mắt vài lần, rồi lặng lẽ bỏ cành cây cũ xuống, lấy một nhánh nhỏ hơn.

Nhánh cây này nhìn bên ngoài trông khá bình dị, nhưng thực sự rất tươi tốt. Này, đến cả quỷ như ta cũng muốn có lộc chứ.

Kuon-kun khẽ cười: ``Vậy mới đúng.''

Ta lườm cậu ta, nhưng không phản bác.

Chẳng biết có linh nghiệm không, nhưng khi cầm nhánh cây trên tay, ta cảm thấy lòng mình nhẹ nhõm hơn một chút.

\dots Thật sự không biết liệu ta có thật sự là succubus hay không nữa.

\ct

Trên đường về, Kuon-kun và cô em gái của anh ấy vẫn lườm ta suốt cả buổi vì chuyện phát tà khí lúc nãy.

Còn ta thì lặng lẽ nghĩ: ``\textit{Hừm\dots\ nếu succubus không có tác dụng với thần thánh, vậy có nghĩa là\dots\ mình hoàn toàn bất lực với họ sao?}''

Ta nhíu mày. Đây là một phát hiện quan trọng!

Có lẽ, ta cần phải tìm hiểu thêm về điều này\dots\ nhưng chắc chắn không phải trong một ngôi chùa.

\stpt{-Histoire 5-}
\stcap{Cuộc gọi mờ ám}

Ta vốn chỉ định làm một mẻ bánh quy đơn giản để ăn vặt ngày Tết, nhưng hóa ra lại vô tình khám phá ra một bí mật mờ ám của Kuon-kun.

Mọi chuyện bắt đầu vào chiều mùng Sáu. Trong lúc ta vẫn còn bận rộn trong bếp, nhà Hikawa yên ắng một cách lạ thường. Ryuuhou thì đang nằm dài trên sô-pha xem TV, còn Kuon-kun thì không thấy bóng dáng đâu.

Ta vừa nhào bột vừa lẩm bẩm: ``\textit{Hừm, bình thường giờ này anh ấy đã ra phòng khách rồi mà\dots}''

Nhưng thôi kệ. Ta có nhiệm vụ quan trọng hơn: làm một mẻ bánh quy thật ngon để chứng minh tài nghệ của mình!

\ct

Sau khoảng một tiếng đồng hồ, bánh quy cuối cùng cũng ra lò.

``\textbf{Xong rồi!!}''

Ta hí hửng xếp từng chiếc bánh nóng hổi lên đĩa, cẩn thận để không làm vỡ. Bánh quy lần này có cho thêm chút mật ong, chắc chắn sẽ cực kỳ ngon.

Ta vừa nhấc đĩa bánh lên thì bất chợt\dots

Giọng của ai đó vọng ra từ tầng dưới.

``\dots Anh không quan tâm các em xây dựng hình tượng idol như thế nào. Đó là lựa chọn của các em. Nhưng điều anh muốn là các em phải xây dựng nó sao cho-''

Hả?

Tôi lập tức căng tai lên nghe tiếp, nhưng đúng lúc đó, Ryuuhou đổi kênh ti-vi, tiếng nhạc ầm ĩ lấn át đi câu thoại tiếp theo của anh ấy.

Không thể nào! Anh ta đang nói chuyện với ai đó về idol?

Là ai nhỉ?

Dù sao thì\dots\ giọng của Kuon-kun nghe có vẻ nghiêm túc hơn bình thường. Điều đó càng làm ta tò mò hơn! Làm sao trò chuyện về idol mà thiếu ta được!

Ta quyết định rồi.

Phải xuống xem chuyện gì đang xảy ra!

Vậy là, với đôi tay cầm đĩa bánh quy, ta phóng như bay từ căn bếp xuống tầng năm. Không thể chần chừ thêm, ta lớn giọng, vang cả khắp hành lang: ``\textbf{Kuon-kun\~{}!! Bánh quy của em làm xong rồi đây!! Anh có muốn ăn thử không!?}''

Ngay khi ta đẩy cửa phòng ra, ta thấy Kuon-kun đang ngồi trước màn hình máy tính, ngoái đầu lại với ánh mắt đầy hoảng hốt.

``Anh đang nói chuyện với ai vậy, Kuon-kun? Để em xem thử-'' ta nói, định đặt khay bánh quy xuống và tiến tới xem, thì\dots

``K-Không có!'' Kuon-kun quay mặt về màn hình, nói một cách đầy hối thúc, ``Thôi được, mấy bữa nữa anh đi làm thì ta kể tiếp nha! Bảo trọng!''

Anh ấy vội đóng cái máy tính xuống\dots

*CỘP!*

Và thế là, cuộc gọi kết thúc. Một cách rất đường đột.

Ta chớp mắt. Ơ\dots\ ủa? Hình như ta vừa phá hỏng cái gì đó?

Bên đầu dây bên kia, chắc hẳn những người kia đang rất khó hiểu. Ta cũng thế, tự dưng lại kết thúc cuộc gọi làm gì vậy?

Anh chủ nhà im lặng vài giây, rồi quay sang ta, thở dài: ``\dots Em có thể đừng hét to thế không?''

Ta nhíu mày, chớp mắt ngây thơ: ``Ơ? Nhưng em đã làm bánh quy cho anh mà? Em muốn được anh nghiệm thử thành quả do chính tay em làm!''

Lúc này, Ryuuhou cũng liếc nhìn bọn ta, khoanh tay nhìn anh trai của cô đầy nghi ngờ. ``\dots Anh vừa gọi cho ai thế?''

Anh ấy im lặng một lúc, ánh mắt liếc sang ta. Có vẻ như anh không muốn nói khi ta vẫn ở đây.

Chà chà, đáng ngờ thật đấy! Một bí mật gì đó mà ta không có quyền được biết sao?

Nhưng thay vì trả lời, Kuon-kun chỉ thản nhiên cầm một chiếc bánh quy lên cắn một miếng, rồi gật đầu: ``Không có gì, ăn bánh thôi.''

Ta chưa kịp phản ứng thì Ryuuhou đã cầm một cái bánh lên, cắn thử. ``Hừm\dots\ cũng ngon đấy.''

Mắt ta sáng rực với lời khen bất ngờ đó: ``Thật á!?''

Và, chính lời nhận xét ngắn gọn đó đã làm ta quên béng chuyện vừa rồi. Ta hào hứng ngồi xuống ăn bánh cùng họ, thưởng thức thành phẩm do ta làm ra.

Ryuuhou là thế, vậy còn Kuon-kun thì sao?

Anh ấy chỉ lặng lẽ nhâm nhi bánh quy và uống trà, như thể chẳng có chuyện gì xảy ra. Nhưng ta biết chắc chắn - rất chắc chắn - rằng cuộc gọi đó có vấn đề.

Và ta sẽ tìm hiểu nó vào một ngày khác! Không sớm thì muộn, anh cũng sẽ phải mở mồm ra mà khai hết thôi, Kuon-kun à!

\end{document}